\documentclass[10.5pt]{article}

\usepackage{amsmath,amssymb}
\usepackage{bm}
\usepackage{geometry}
\usepackage{ctex}
\usepackage{setspace}
\usepackage{graphicx}
\usepackage{float}

\geometry{a4paper,margin=2cm}

%========================
% 排版设置
%========================

\setstretch{1.2}
\setlength{\parskip}{4pt}
\setlength{\parindent}{2em}

\setlength{\abovedisplayskip}{5pt}
\setlength{\belowdisplayskip}{5pt}
\setlength{\abovedisplayshortskip}{3pt}
\setlength{\belowdisplayshortskip}{3pt}

\setcounter{secnumdepth}{3}


\title{\textbf{非互易输运：理论机制与输运特征}}

\author{Super}

\date{}


\begin{document}

\maketitle


\section{非互易输运定义}

在传统线性输运体系中，电流与电压满足Onsager互易关系：

\begin{equation}
V(I)=-V(-I)
\end{equation}

因此电阻满足：

\begin{equation}
R(+I)=R(-I)
\end{equation}


然而，在某些具有特殊对称性破缺的量子材料中，
正、反方向输运响应不再等价：

\begin{equation}
R(+I)\neq R(-I)
\end{equation}

这种现象称为非互易输运
(Nonreciprocal transport)。


非互易输运本质上来源于体系中的二阶非线性响应：

\begin{equation}
J=\sigma E+\chi E^2+\cdots
\end{equation}


其中：

\begin{itemize}
\item $J$为电流密度；
\item $E$为外加电场；
\item $\sigma$为线性电导；
\item $\chi$为二阶非线性输运系数。
\end{itemize}


其中：

\begin{equation}
J^{(2)}=\chi E^2
\end{equation}

对应非互易输运贡献。



\section{非互易输运理论机制}


\subsection{对称性破缺条件}

非互易输运受到体系对称性的限制。

根据Onsager关系，
若体系同时满足时间反演和空间反演对称性，
则二阶响应必须为零：

\begin{equation}
\chi=0
\end{equation}


因此非互易输运通常要求：

\begin{equation}
\boxed{
\text{时间反演对称性破缺}
+
\text{空间反演对称性破缺}
}
\end{equation}


时间反演操作：

\begin{equation}
t\rightarrow -t
\end{equation}


使电子动量和自旋发生反转：

\begin{equation}
\bm{k}\rightarrow-\bm{k}
\end{equation}


外加磁场：

\begin{equation}
B\neq0
\end{equation}

或磁有序：

\begin{equation}
M\neq0
\end{equation}

均可以破坏时间反演对称性。


另一方面，
空间反演破缺：

\begin{equation}
\bm{r}\rightarrow-\bm{r}
\end{equation}

可以使电子态具有方向相关性，
从而允许非互易响应。



\subsection{Rashba自旋轨道耦合机制}

在界面体系或非中心对称晶体中，
结构电场能够诱导Rashba自旋轨道耦合：

\begin{equation}
H_R=
\alpha_R
(\bm{k}\times\bm{\sigma})
\cdot
\bm{E}
\end{equation}


其中：

\begin{itemize}
\item $\alpha_R$为Rashba耦合系数；
\item $\bm{k}$为电子动量；
\item $\bm{\sigma}$为电子自旋。
\end{itemize}


Rashba效应导致：

\begin{equation}
E(k)\neq E(-k)
\end{equation}


使电子沿不同方向运动时具有不同散射概率，
产生非互易输运。



\section{非互易输运理论模型}


非互易电阻通常定义为：

\begin{equation}
\Delta R=
\frac{R(+I)-R(-I)}{2}
\end{equation}


也可以定义归一化非互易系数：

\begin{equation}
\gamma=
\frac{R(+I)-R(-I)}
{R(+I)+R(-I)}
\end{equation}


在磁场作用下，
非互易响应满足：

\begin{equation}
\Delta R(B)
=
-\Delta R(-B)
\end{equation}


即磁场方向改变时，
非互易信号发生符号反转。


通常：

\begin{equation}
\Delta R\propto B
\end{equation}

表明其与磁场诱导的时间反演破缺有关。



\section{非互易输运测试方法与实验特征}


\subsection{直流电流反转测量}


直流电流反转是最直接的非互易输运测试方法。


实验中分别施加：

\begin{equation}
+I,\quad -I
\end{equation}


并测量：

\begin{equation}
V(+I),\quad V(-I)
\end{equation}


对于普通欧姆输运：

\begin{equation}
V(+I)=-V(-I)
\end{equation}


而非互易体系满足：

\begin{equation}
V(+I)+V(-I)\neq0
\end{equation}


定义非互易电压：

\begin{equation}
\Delta V=
\frac{V(+I)+V(-I)}{2}
\end{equation}


该方法可以直接观察输运方向不对称。



\subsection{二次谐波输运测量}


交流电流：

\begin{equation}
I=I_0\sin(\omega t)
\end{equation}


电压响应展开为：

\begin{equation}
V=
R_1I+R_2I^2+\cdots
\end{equation}


其中：

\begin{equation}
R_1I
\end{equation}

对应线性输运。


二阶项：

\begin{equation}
R_2I^2
\end{equation}

产生二次谐波信号：

\begin{equation}
V_{2\omega}\propto I^2
\end{equation}


因此：

\begin{equation}
V_{2\omega}\neq0
\end{equation}

是非互易输运的重要实验判据。



\subsection{磁场依赖测量}


由于非互易输运通常需要时间反演破缺，
外加磁场是重要调控参数。


实验测量：

\begin{equation}
\Delta R(B)
\end{equation}


其典型特征：

\begin{equation}
\Delta R(-B)
=
-\Delta R(B)
\end{equation}


通过磁场反转实验，
可以判断信号是否来源于磁性相关机制。



\subsection{温度依赖测量}


非互易输运通常在电子相变附近增强。


实验中测量：

\begin{equation}
\gamma(T)
\end{equation}


或者：

\begin{equation}
V_{2\omega}(T)
\end{equation}


典型增强区域包括：

\begin{itemize}
\item 超导转变温度$T_c$附近；
\item 磁有序温度$T_N$附近；
\item 金属-绝缘体转变温度$T_{MIT}$附近。
\end{itemize}


增强原因包括：

\begin{itemize}
\item 磁涨落增强；
\item 电子关联增强；
\item 对称性破缺态形成。
\end{itemize}



\subsection{角度依赖测量}


对于具有晶体各向异性的材料，
旋转磁场方向可以研究非互易输运的空间对称性。


测量：

\begin{equation}
\Delta R(\theta)
\end{equation}


其中$\theta$为磁场方向与晶体轴夹角。


角度依赖可以揭示：

\begin{itemize}
\item Rashba自旋轨道耦合；
\item 磁结构方向；
\item 晶体对称性破缺。
\end{itemize}



\subsection{Hall输运联合测量}


对于强关联材料，
非互易输运通常伴随载流子重构。


Hall系数：

\begin{equation}
n_H=
\frac{1}{eR_H}
\end{equation}


结合：

\begin{equation}
\rho(T),
\quad
n_H(T),
\quad
V_{2\omega}(T)
\end{equation}


可以分析：

\begin{itemize}
\item 载流子浓度变化；
\item 电荷序转变；
\item 磁性散射变化。
\end{itemize}



\section{强关联氧化物中的非互易输运}


强关联氧化物中，
电子关联、磁序以及晶格畸变可以共同产生非互易输运。


以稀土镍酸盐：

\begin{equation}
RNiO_3
\end{equation}

为例，

NdNiO$_3$具有：

\begin{itemize}
\item 金属-绝缘体转变；
\item 电荷序；
\item 磁有序；
\item 强电子关联。
\end{itemize}


低温区域，
局域磁矩、自旋涨落以及结构反演破缺可能导致：

\begin{equation}
R(+I)\neq R(-I)
\end{equation}


因此，通过：

\begin{equation}
V_{2\omega}-T
\end{equation}

以及：

\begin{equation}
\Delta R-B-T
\end{equation}


可以研究镍酸盐体系中的隐藏对称性破缺和电子态演化。



\section{总结}


非互易输运是一种由时间反演对称性破缺和空间反演对称性破缺共同导致的二阶非线性输运现象。


主要实验特征包括：

\begin{enumerate}

\item 正反电流电阻不等：

\[
R(+I)\neq R(-I)
\]


\item 二次谐波响应：

\[
V_{2\omega}\propto I^2
\]


\item 磁场反转信号反转：

\[
\Delta R(-B)
=
-\Delta R(B)
\]


\item 在电子相变附近出现增强。

\end{enumerate}


非互易输运为研究强关联氧化物、
磁性体系以及拓扑量子材料中的隐藏电子态提供了重要实验方法。


\end{document}